\section{Main Result}

\begin{definition}[Classical target class]
The classical target class consists of classical solution trajectories of the original 3D incompressible Navier--Stokes equation on \([0,T)\times\mathbb{R}^3\), generated by smooth compactly supported divergence-free initial data, and for which each theorem-critical control mechanism is derived on the same equation surface with no essential dependence on modified-equation terms.
\end{definition}

\begin{theorem}[Conditional source-grounded route theorem]
\label{thm:classical-closure}
Assume the four bridge packages from this manuscript hold on one admissibility
surface for the classical target class. Then the tracked
scale/monotone/compactness/gradient singularity pathway does not occur on
\([0,T)\). In particular, the dangerous high-frequency cascade is suppressed,
the monotone functional \(Q(t)\) remains explicit as same-surface global
control infrastructure, the nonlinear term closes on one classical
approximation surface with no defect, and the fourth bridge controls the
derivative-scale channel that would otherwise reopen singularity.
\end{theorem}

The lane also carries a separate NS-only projected-flow track, which remains
secondary and is not a substitute for the source-grounded four-bridge route.
Its formal closure package is:

\[
\begin{aligned}
&\text{D.1 projected NC-flow well-posedness}
\to \text{D.2 global coercive energy estimate}\\
&\to \text{D.3 exact commutative shadow descent}
\to \text{D.4 continuation from bounded ontic control}\\
&\to \text{D.5 final Clay integration.}
\end{aligned}
\]

The historically hard point in that sidecar is D.2. The closure hypothesis required there is
the coherent/remainder split
\[
\mathbb P_DN_D(X)=N_{\mathrm{coh}}(X)+N_{\mathrm{rem}}(X),
\]
together with the exact carrier bound
\[
\begin{aligned}
&\|[\nabla_{D,\omega},N_{\mathrm{coh}}(X)]\|
+\|N_{\mathrm{rem}}(X)\|\\
&\qquad
+\|[\nabla_{D,\omega},N_{\mathrm{rem}}(X)]\|\\
&\le
\varepsilon\Big(
\|\Delta_{D,\omega}X\|
+\|\Delta_{D,\omega}^{1/2}\Omega_D(X)\|
\Big)
+C_\varepsilon\Psi(\mathcal E_D(X)).
\end{aligned}
\]
where
\[
\mathcal E_D(X)=\|\Delta_{D,\omega}^{1/2}X\|^2+\|\Omega_D(X)\|^2.
\]
The exact lift needed downstream is the strict-shadow package
\[
\begin{aligned}
\mathcal C(\partial_t X)&=\partial_t u,\qquad
\mathcal C(\Delta_DX)=\Delta u,\\
\mathcal C(\mathbb P_DX)&=\mathbb P_{\mathrm{Leray}}u,\qquad
\mathcal C(\mathbb P_DN_D(X))=\mathbb P_{\mathrm{Leray}}(u\cdot\nabla u),
\end{aligned}
\]
with the continuation bridge
\[
\|u\|_{H^s}\le C\|X\|_{H_D^s},\qquad s>\frac52.
\]

The route theorem is organized around four bridge propositions on one declared
surface. The first is a scale barrier against singular cascade transfer. The
second is a monotone functional that controls the global energy/enstrophy
layer. The third is a nonlinear-compactness package that closes the convective
term. The fourth is the source-faithful curvature-based bridge together with
its feedback into the scale barrier. The theorem above is conditional on those
four bridge packages being available on one admissibility surface for the
classical target class.

The realized classical carrier and exact-pullback layer are now explicit in the
realized-carrier note, and the external validation of the imported
Leray/Stokes/Aubin-Lions package is recorded in the companion external
classical-validation note. In
particular, on \(\mathbb R^3\) the compactness bridge is read in the standard
way: Aubin--Lions--Simon gives strong convergence on bounded balls, and the
global \(L^2\) conclusion used by the route comes only from the admissible
energy-tail bound already assumed on the same theorem surface
\cite{Temam2001,Simon1987,BoyerFabrie2013}. No new theorem-local sidecar lemma
is claimed there. Bridge IV is stated separately as a pure Euclidean
gradient-control theorem in \texttt{sections/bridge-gradient-control.tex},
while the remaining open surface is the closed-loop warrant on the declared
theorem surface.
